% Supplementary Note 1: Study 1's hypotheses 3 and 4. Numbers through macros only.
\suppnote{Study 1: hypotheses 3 and 4}

This note reports in full Study~1's third and fourth hypotheses (H3 and H4). The main text's Results on Study~1 give only their verdicts, in one sentence.

Study~1's \SOneN{} participants rated ten AI-made images, five showing a human or humanoid figure and five not, for creativity and for effort on 1-to-5 scales. A label told each participant that the artist, Alex, had made the images either with design and sketching tools (the Human label; \SOneNHuman{} participants) or by prompting an AI system (the AI label; \SOneNAI{} participants). Each participant saw one label and rated every image. We call the mean rating under the Human label minus the mean under the AI label the label's penalty; a positive penalty means that the AI label lowered the rating. Each estimate comes from a linear mixed model, a regression that gives each participant and each image its own baseline, fitted to each rating separately. Study~1 was registered on the Open Science Framework (OSF), an online archive for study plans, on 8 October 2026, after its data had been collected and analyzed. H3's bound and decision table, and the quantity that decides H4, were therefore set with the estimates already known. Supplementary Table~\ref{tab:supp-verdicts} lists both verdicts beside the other registered tests.

\paragraph*{H3: the same penalty for images with and without a human figure.}
H3 predicted that the penalty would not differ between the two kinds of image. It is a registered hypothesis decided outside Study~1's set of confirmatory tests, by its own rule rather than by a corrected $p$-value: a significance test read with an equivalence test. The significance test asks whether the difference in penalty between the two kinds of image is distinguishable from zero. Its $p$-value is the chance of a difference at least this large if there were none, and a difference is called significant when $p$ is below .05. An equivalence test asks whether an effect is shown to lie inside a stated bound. Judged alone, as H3's is, it passes when the effect's 90\% confidence interval (CI), the interval that matches this test, lies wholly inside the bound, which shows that the effect is small, not that it is zero.\cite{schuirmann1987comparison,lakens2017equivalence} H3's bound was $\pm$\SOneHThreeBound{} points. Read together, the two tests decide H3: a difference that is not significant and lies inside the bound supports H3, and one that is not significant but is not shown inside the bound leaves H3 inconclusive, because the design cannot then separate no difference from a small one.

On effort, H3 was supported: the difference in penalty was not significant ($p$~\SOneHThreeEffP) and lay inside the bound (\SOneHThreeEffEst{} points; 90\% CI \SOneHThreeEffLoNinety{} to \SOneHThreeEffHiNinety; equivalence $p$~\SOneHThreeEffPTost). The bound was set after the analyses, so this pass is weaker evidence than a pass against a bound fixed in advance. On creativity, H3 was inconclusive (\SOneHThreeCreatEst{} points; 90\% CI \SOneHThreeCreatLoNinety{} to \SOneHThreeCreatHiNinety; $p$~\SOneHThreeCreatP), because the interval crossed the bound. The equivalence test's power, its chance of passing if the true difference were zero, was only \SOneHThreeCreatPower, so that failure says little. By the registered rule, a failed equivalence test is uninformative whenever this power is below 80\%, since the failure then reflects the width of the interval this design produces, not the difference itself. Power on effort was much the same (\SOneHThreeEffPower), but a pass stands whatever the power, because the equivalence test limits the chance of wrongly declaring a difference small. Were the true difference a tenth of a point, the test would pass with probability \SOneHThreeCreatPowerSmall{} on creativity and \SOneHThreeEffPowerSmall{} on effort. These power figures come from simulated data and describe the design, not the observed data. They are a registered description: a result that the registration asks us to report, with no test.

Both verdicts rest on the primary model. Each verdict was also to hold in a fuller model, one that further lets the label's effect vary from image to image and the effect of image kind vary from person to person, unless that model was singular. A singular fit estimates a variance or correlation at the edge of its possible range, a sign that the data cannot support that many quantities. The fuller model was singular on both ratings, so under the registered fallback the verdicts rest on the primary model alone. The singular fuller model gave the same readings (90\% CI \SOneHThreeEffMaxLoNinety{} to \SOneHThreeEffMaxHiNinety{} on effort, \SOneHThreeCreatMaxLoNinety{} to \SOneHThreeCreatMaxHiNinety{} on creativity).

\paragraph*{H4: more spread between people under the AI label.}
H4 predicted that, under the AI label, ratings would spread out more between people and less across one person's ratings of different images. It is a registered confirmatory hypothesis: Study~1's registration names it among the study's confirmatory tests. Its registered criterion, on each rating, was an AI-to-Human ratio of the variance (the spread) between people with a 95\% CI wholly above 1, and a ratio within people with a 95\% CI wholly below 1. A 95\% CI is the range of values the data do not rule out at the 95\% level; here each came from recomputing the ratio after drawing participants again at random, 1,000 times, within each label group. The criterion was met on effort: the ratio was \SOneHFourEffRatioBetween{} between people (\SOneHFourEffRatioBetweenLo{} to \SOneHFourEffRatioBetweenHi) and \SOneHFourEffRatioResidual{} within people (\SOneHFourEffRatioResidualLo{} to \SOneHFourEffRatioResidualHi). It was not met on creativity. There the ratio between people was \SOneHFourCreatRatioBetween{} (\SOneHFourCreatRatioBetweenLo{} to \SOneHFourCreatRatioBetweenHi), but the ratio within people, \SOneHFourCreatRatioResidual, had an interval that reached \SOneHFourCreatRatioResidualHi. Because H4 required both ratings, it was not supported as registered.

Study~1's four confirmatory tests were judged together after Holm's correction,\cite{holm1979simple} which raises each $p$-value in a set of tests so that the chance of calling any result in the set significant by mistake stays at .05. H4's row in that set survived the correction (Holm $p$~\SOneHFourPHolm), but that row does not decide H4. It is only the likelihood-ratio test of the spread between people on creativity, the rating with the weaker result, which the registration placed in the set; the test compares a model in which that spread differs by label with one in which it does not. Passing it is not passing H4. Three further points limit the reading. The two ratios come from fits with different structures, and the fit for the spread within people gives images no baseline of their own, so that spread also carries differences between images. In simulated data with ratios equal to those observed, the test of the spread between people would reject in \SOneHFourCreatResolutionPowerPct{} of samples on creativity and \SOneHFourEffResolutionPowerPct{} on effort, so the creativity half was the harder to detect (registered description of what the design can resolve). And the quantity that decides H4 was chosen after the analyses had been run.

As description that the registration places outside the H4 criterion, more AI-label participants gave the same rating to seven or more of the ten images. On effort, \SOneHFourEffSameRatingAIN{} of \SOneNAI{} AI-label participants did so, against \SOneHFourEffSameRatingHumanN{} of \SOneNHuman{} Human-label participants; on creativity, \SOneHFourCreatSameRatingAIN{} against \SOneHFourCreatSameRatingHumanN. These counts describe the ratings and test nothing. Supplementary Fig.~\ref{figS:raincloud} shows each participant's mean rating under each label; that these means spread more under the AI label is likewise description, not a test of H4.

\begin{figure}[tbp]
\centering
\makebox[\textwidth][c]{\includegraphics[width=183mm]{figS_study1_raincloud}}
\caption{\textbf{Study~1: each participant's mean rating under each label.} Each small dot is one participant's mean rating of the ten originals, filled under the Human label (\SOneNHuman{} participants) and open under the AI label (\SOneNAI{}), stacked where values repeat. The outline above the dots, a half-violin, is a smoothed curve of how they spread. The large marks are the two label means, which here equal the means predicted by Study~1's statistical model. The right column is the gap between them, the Human-label mean minus the AI-label mean, with its 95\% confidence interval, as in main-text Fig.~2a; these gaps are the registered confirmatory tests of H1, that the AI label lowers both ratings. The two groups' dots overlap, so a gap is a difference between group means, not a difference that every participant shows. The dots also spread more under the AI label, more on effort than on creativity; this spread is descriptive. H4, the registered confirmatory test that the AI label spreads ratings more between people and less within each person, met its criterion on effort only, so it was not supported as registered.}
\label{figS:raincloud}
\end{figure}
